Las pruebas de carga se ejecutan con Locust contra el \textit{endpoint}
\texttt{/api/v1/} para verificar el \cref{rnf:concurrency}. Se han
ejecutado tres escenarios con perfiles de carga representativos:

\begin{figure}[Fase 2: percentiles de latencia]%
              {fig:cdetail-phase2-latency}%
              {Fase 2 (corrección concurrente): percentiles de latencia
               por \textit{endpoint}.}
  \image{}{}{load_test/graphs/phase2/01_latency_percentiles}
\end{figure}

\begin{enumerate}

  \item \textbf{Ingesta sostenida.} Simulación de una cabina de escaneo
  entregando páginas a través de la \acs{api} de ingesta. El sistema
  procesó 14\,260 peticiones a una tasa de 23,8 \textit{RPS} sin
  errores, con latencia mediana inferior a un milisegundo.

  \item \textbf{Corrección concurrente.} Simulación de correctores
  anotando y calificando simultáneamente. El sistema procesó cerca de
  5\,900 peticiones sin errores, con latencias medianas de 14 a 44
  milisegundos por operación, satisfaciendo el objetivo de p95 inferior
  a 500 milisegundos para los \textit{endpoints} \acs{crud} (\cref{fig:cdetail-phase2-latency}).

  \item \textbf{Pico de publicación.} Simulación de la situación más
  exigente: la publicación de calificaciones disparando un acceso masivo
  de alumnos consultando y descargando sus exámenes. Bajo este perfil,
  el sistema mantiene la latencia de los \textit{endpoints} ligeros pero
  presenta degradación significativa en
  \texttt{student/instance/download} (alrededor del 21\,\% de errores
  \acs{http} 500 y latencias medianas de 24 segundos), lo que excede el
  objetivo de respuesta sin errores. Esto es debido a la gran carga de trabajo
  que implica el endpoint (conexión con redis y la base de datos,
  recuperación del pdf y generación de un nuevo pdf con la marca de agua en todas
  las páginas).

\end{enumerate}

Los dos primeros escenarios validan el cumplimiento del
\cref{rnf:concurrency} en operación normal. El tercer escenario
identifica un cuello de botella consistente con el uso del servidor de
desarrollo de Django (\texttt{runserver}) en lugar de un servidor
\acs{wsgi} dedicado, lo que se aborda en la línea de trabajo futuro
descrita en el \cref{CAP:CONCLUSIONS}. El detalle completo de los
escenarios, los gráficos de latencia y la lista de errores observados
se incluyen en \cref{AP:TEST-RESULTS}.
